\documentclass{article}
\usepackage[hmargin={0.5in, 0.5in}, vmargin={0.8in, 0.8in}]{geometry}
\usepackage{booktabs}
\usepackage[table]{xcolor}
\usepackage{hyperref}
\hypersetup{
  colorlinks=true,
  linkcolor=blue,
  urlcolor=blue
}
\usepackage[anythingbreaks]{breakurl}
\usepackage{multicol}
\usepackage[compact]{titlesec}

\newcommand{\pkg}[1]{{\normalfont\fontseries{b}\selectfont #1}}
\let\proglang=\textsf
\let\code=\texttt

\usepackage{enumitem}
\setlist{parsep=0pt, leftmargin=4mm, topsep=0pt, itemsep=5pt}

\usepackage{fancyhdr}
\pagestyle{fancy}
\renewcommand{\headrulewidth}{0.4pt}
\renewcommand{\footrulewidth}{0.4pt}

\lhead{\sf STAT 5255/3255}
\rhead{Fall 2026}
\lfoot{\sf jun.yan@uconn.edu}
\rfoot{\url{https://statcomp.org/}}

\begin{document}

\title{STAT 5255/3255: Introduction to Data Science}
\date{August 31, 2026}
\maketitle

\thispagestyle{fancy}

\begin{multicols}{2}
\begin{description}
\item[Instructor:]
  Jun Yan\\
  AUST 328\\
  jun.yan@uconn.edu

\item[Lectures:]
  MoWe 8:40AM--9:55AM\\
  AUST 103

\item[Office Hours:]
  MoTh: 12:00--1:00 pm
  or by appointment.

\item[Email Etiquette:]
  Learn the basics of professional email etiquette when writing to
  professors:
  \href{https://marktomforde.com/academic/undergraduates/Email-Etiquette.html}{Guidelines for Writing to Your Professors}.
  Use email rather than HuskyCT messages to contact the instructor.

\item[Prerequisites:] \hspace{0pt}

  STAT 5255: Open to graduate students in Statistics and to others with
  permission. Not open for credit to students who have passed STAT 3255.

  STAT 3255: STAT 2255 and STAT 3115Q (or 3215Q), or instructor consent.

  Prior experience with Python and regression is assumed.

\item[Course Description:]
  This course introduces data science methods for storing, processing, and
  analyzing data effectively and drawing inferences. Topics
  include project management, data preparation, data visualization,
  statistical models, machine learning, distributed computing, and ethics.

\item[Course Materials:]
  Lecture notes, assignments, sample code, datasets, and other private course
  materials will be posted on HuskyCT at \url{https://lms.uconn.edu/}.
  Students are responsible for checking HuskyCT regularly.

\item[Recommended Texts:]\hspace{0pt}
  \begin{enumerate}
  \item
    ``Python Data Science Handbook: Essential Tools for Working with Data,''
    First Edition, by Jake VanderPlas, O'Reilly Media, 2016.
    \url{https://jakevdp.github.io/PythonDataScienceHandbook/}.
  \item
    ``Python for Data Analysis: Data Wrangling with Pandas, NumPy, and
    IPython,'' Third Edition, by Wes McKinney, O'Reilly Media, 2022.
    \url{https://wesmckinney.com/book/}.
  \item
    ``Veridical Data Science: The Practice of Responsible Data Analysis and
    Decision Making,'' by Bin Yu and Rebecca L. Barter, MIT Press, 2024.
    \url{https://vdsbook.com}.
  \item
    \href{https://www.practicaldatascience.org}{``Practical Data Science in
    Python''} by Kyle Bradbury and Nick Eubank.
  \end{enumerate}

\item[Computing:]
  \proglang{Python} will be the primary computing environment. Students may
  choose an interactive development environment such as Jupyter, VS Code, or
  Emacs. Follow the
  \href{https://peps.python.org/pep-0008/}{Python Code Style Guide}.

  \proglang{Git} will be used for project management and collaboration, and
  Classroom 50, an open-source replacement for the retired GitHub Classroom,
  will be used for homework submission and grading. Practice
  Git with the \href{https://gitexercises.fracz.com}{online exercises}.

  Command-line operation is essential for reproducibility and automation.
  Learn the basics from a suitable tutorial, such as the
  \href{https://ubuntu.com/tutorials/command-line-for-beginners}{Ubuntu
  command-line tutorial}.

  \proglang{Quarto} will be used for student presentation contributions,
  homework assignments, and projects. Always be ready to learn new tools.

\item[Class Notes:]
  The course uses two collections of notes. The instructor-authored
  \href{https://statds.github.io/ids-book/}{Introduction to Data Science book}
  contains the permanent course notes. The
  \href{https://statds.github.io/ids-f26/}{Fall 2026 class notes}
  organize and link the student-created RevealJS presentations whose sources
  have been accepted into the
  \href{https://github.com/statds/ids-f26}{Fall 2026 GitHub repository}.

  Past class notes are available for
  \href{https://statds.github.io/ids-f25}{Fall 2025} and
  \href{https://statds.github.io/ids-s26}{Spring 2026}.

\item[Fall 2026 Calendar:]
  The semester begins Monday, August 31. Classes do not meet on Labor Day,
  Monday, September 7. Thanksgiving recess runs Sunday, November 22, through
  Sunday, November 29. The final two instructional weeks, Monday, November 30,
  through Friday, December 11, are reserved for final-project presentations.
  Friday, December 11, is the last day of classes. Final examinations run
  Monday, December 14, through Sunday, December 20.

\item[The First Two Weeks:]
  The first two weeks are intentionally demanding because students must learn
  the computing and submission workflow under real deadlines. The first two
  assignments are therefore substantial, and by the end of the second week
  every student is expected to be able to set up the required environment,
  complete the work, and submit it correctly through Classroom 50. Seek
  help immediately when a setup or submission problem arises; not knowing how
  to submit work by midsemester is not acceptable.

  Use these two weeks to make an honest decision about whether you have the
  time and commitment required for the course. Students who cannot meet these
  early expectations should drop the course by the applicable deadline. Do
  not assume that this warning applies only to other students.

\item[Grading:]
The course grade will be based on the following components:
\begin{center}
  \rowcolors{2}{blue!20}{gray!20}
  \begin{tabular}{lrr}
    \toprule
    Category              & 3255  & 5255\\
    \midrule
    Style                 & 5\%  & 5\%\\
    Participation         & 5\%  & 5\%\\
    Homework              & 20\% & 20\%\\
    Topic Presentations   & 20\% & 20\%\\
    Midterm Project       & 20\% & 20\%\\
    Final Presentation    & 30\% & 0\%\\
    Final Report          & 0\%  & 30\%\\
    \bottomrule
  \end{tabular}
\end{center}

\item[Style: 5 points]
  Homework and midterm submissions should clearly identify the problem being
  addressed or quote the question. A submission should present a clear train
  of thought. One style point may be deducted for each occurrence of
  illegible, ambiguous, incoherent, or explicitly nonconforming work until the
  five style points are exhausted.

\item[Participation: 5 points]
  Attendance and participation during final presentations from November 30
  through December 11 are required. Missing this part of the course forfeits
  the five participation points. Each student is expected to ask at least two
  substantive questions during the presentation period.

\item[Homework: 20 points]\hspace{0pt}
\begin{itemize}
\item Homework will be assigned roughly weekly during the first two-thirds of
  the semester. Students may consult one another or the instructor, but each
  student must submit independent work.
\item Project-completion instructions must be followed closely.
\item Work is due at the stated date and time unless additional flexibility is
  announced. Work submitted within 48 hours after the deadline receives a
  linearly decreasing maximum score from 95\% to 50\%. Work submitted more
  than 48 hours late will not be graded.
\item Submitted work that duplicates another solution receives no credit.
\end{itemize}

\item[Topic Presentations: 20 points]
  Each student selects a topic from the public task board, obtains approval,
  prepares a RevealJS slide deck, and presents it to the class. The
  presentation must be submitted for preview at least one week before the
  scheduled presentation; approval and timely preview account for 4 points.
  The oral presentation accounts for 10 points. Successful completion of a pull request containing
  the revised slide source and topical catalog entry within two weeks after
  the presentation accounts for the remaining 6 points.

\item[Midterm Project: 20 points]
  The midterm project will be assigned and will require completion of the
  analysis component of a full data-science project. Details, dates, and
  deliverables will be posted on HuskyCT.

\item[Final Project: 30 points]
  The final-project topic is chosen by the student and should demonstrate the
  complete cycle of a real data-science project. Data-science competition
  sites such as \url{https://kaggle.com} may provide suitable ideas.

\begin{itemize}
\item Undergraduate students in STAT 3255 will give 15-minute presentations during
  the final two instructional weeks, November 30--December 11. The order will
  be derived from the random order determined by the class.

\item Graduate students in STAT 5255 will submit a final written report by
  11:59 pm on Friday, December 11, 2026.

\item A one-page proposal is due by 11:59 pm on Friday, October 23. The proposal
  accounts for 5 of the 30 final-project points. A proposal submitted during
  the following 24-hour grace period can receive at most 50\% of the proposal
  credit. A proposal submitted after the grace period receives no proposal
  credit but must still be completed.

\item By 11:59 pm on Friday, October 30, each student must arrange a meeting with
  the instructor about the project, with the meeting to occur by Friday,
  November 6. The meeting accounts for 5 points. Failure to arrange or attend
  the meeting results in loss of the corresponding credit.
\end{itemize}

\item[Generative Artificial Intelligence:]
  Students are encouraged to use AI tools to help set up their computing
  environments, learn data-science operations and workflows, and brainstorm
  ideas. Each student must complete their own homework, topic presentation,
  midterm project, and final project. Any AI use related to submitted work
  must be disclosed. Students remain responsible for verifying
  all code, claims, citations, and results and for being able to explain and
  defend their work.

\item[Avoiding Plagiarism:]
  Students are responsible for avoiding plagiarism and meeting University
  expectations for academic integrity. Consult the UConn Library resources on
  \href{https://guides.lib.uconn.edu/plagiarism}{understanding plagiarism} and
  \href{https://guides.lib.uconn.edu/citing}{citing sources}.

\item[University Policies and Academic Integrity:]
  This course adheres to the UConn
  \href{https://policy.uconn.edu/2023/07/11/academic-scholarly-and-professional-integrity-and-misconduct-aspim-policy-on/}{Academic,
  Scholarly, and Professional Integrity and Misconduct policy}.

\item[Copyright:]
  Copyrighted course materials are provided only for use by students enrolled
  in the course and may not be retained or redistributed except as authorized.
  Public collaborative materials are governed by the license stated in their
  repository.

\item[Recommendation Letters:]
  A useful recommendation letter evaluates attributes such as motivation,
  academic preparation, research or teaching potential, teamwork,
  communication, and maturity. Students seeking a letter should demonstrate
  those attributes through their work and interactions with the instructor.

\item[Disclaimer:]
  The instructor reserves the right to change the syllabus when circumstances
  require it. Changes will be communicated to the class.
\end{description}
\end{multicols}

\end{document}
